\documentclass[11pt]{article}
% The preamble every project's paper shares: % The preamble every project's paper shares: % The preamble every project's paper shares: \input{../../preamble} from
% projects/<project>/paper/. Boards come from chessboard and xskak, plots from
% pgfplots; every number a paper prints comes from its generated.tex.
\usepackage[LSBC4,T1]{fontenc}
\usepackage[utf8]{inputenc}
\usepackage{lmodern}
\usepackage[margin=1in]{geometry}
\usepackage{microtype}
\usepackage{amsmath}
\usepackage{booktabs}
\usepackage{tabularx}
\usepackage{array}
\usepackage{xcolor}
\usepackage[obeyspaces,spaces]{url}
\usepackage{xskak}
\usepackage{pgfplots}
\pgfplotsset{compat=1.18}
\usepackage{caption}
\usepackage{subcaption}
\usepackage{enumitem}
\usepackage[section]{placeins}
\usepackage[hidelinks]{hyperref}

\setchessboard{boardfontencoding=LSBC4, boardfontfamily=skaknew, blackfieldcolor=black!22,
  whitefieldcolor=white, setfontcolors, showmover=false, label=true, labelfontsize=6pt,
  boardfontsize=14pt}

% A move or line in UCI, and a position in FEN, both breakable at a slash.
\newcommand{\uci}[1]{\texttt{#1}}
\newcommand{\fen}[1]{{\footnotesize\path{#1}}}
\newcolumntype{L}{>{\raggedright\arraybackslash}X}

\setlength{\parskip}{0.4em}
\setlength{\parindent}{0pt}
 from
% projects/<project>/paper/. Boards come from chessboard and xskak, plots from
% pgfplots; every number a paper prints comes from its generated.tex.
\usepackage[LSBC4,T1]{fontenc}
\usepackage[utf8]{inputenc}
\usepackage{lmodern}
\usepackage[margin=1in]{geometry}
\usepackage{microtype}
\usepackage{amsmath}
\usepackage{booktabs}
\usepackage{tabularx}
\usepackage{array}
\usepackage{xcolor}
\usepackage[obeyspaces,spaces]{url}
\usepackage{xskak}
\usepackage{pgfplots}
\pgfplotsset{compat=1.18}
\usepackage{caption}
\usepackage{subcaption}
\usepackage{enumitem}
\usepackage[section]{placeins}
\usepackage[hidelinks]{hyperref}

\setchessboard{boardfontencoding=LSBC4, boardfontfamily=skaknew, blackfieldcolor=black!22,
  whitefieldcolor=white, setfontcolors, showmover=false, label=true, labelfontsize=6pt,
  boardfontsize=14pt}

% A move or line in UCI, and a position in FEN, both breakable at a slash.
\newcommand{\uci}[1]{\texttt{#1}}
\newcommand{\fen}[1]{{\footnotesize\path{#1}}}
\newcolumntype{L}{>{\raggedright\arraybackslash}X}

\setlength{\parskip}{0.4em}
\setlength{\parindent}{0pt}
 from
% projects/<project>/paper/. Boards come from chessboard and xskak, plots from
% pgfplots; every number a paper prints comes from its generated.tex.
\usepackage[LSBC4,T1]{fontenc}
\usepackage[utf8]{inputenc}
\usepackage{lmodern}
\usepackage[margin=1in]{geometry}
\usepackage{microtype}
\usepackage{amsmath}
\usepackage{booktabs}
\usepackage{tabularx}
\usepackage{array}
\usepackage{xcolor}
\usepackage[obeyspaces,spaces]{url}
\usepackage{xskak}
\usepackage{pgfplots}
\pgfplotsset{compat=1.18}
\usepackage{caption}
\usepackage{subcaption}
\usepackage{enumitem}
\usepackage[section]{placeins}
\usepackage[hidelinks]{hyperref}

\setchessboard{boardfontencoding=LSBC4, boardfontfamily=skaknew, blackfieldcolor=black!22,
  whitefieldcolor=white, setfontcolors, showmover=false, label=true, labelfontsize=6pt,
  boardfontsize=14pt}

% A move or line in UCI, and a position in FEN, both breakable at a slash.
\newcommand{\uci}[1]{\texttt{#1}}
\newcommand{\fen}[1]{{\footnotesize\path{#1}}}
\newcolumntype{L}{>{\raggedright\arraybackslash}X}

\setlength{\parskip}{0.4em}
\setlength{\parindent}{0pt}

% Written by projects/video-coaching/paper/generate.mjs; do not edit by hand.

\newcommand{\CandidateBullet}{%
\begin{quote}\small
the video coaches the way a strong coach works with a player at his level, in every moment segment, from the evidence alone and inside the video grammar's chapters and rules:
\begin{itemize}[nosep]
\item ask before you tell: the segment's first pause asks what he would play there, open (``What would you play here? Pick two moves.''), before any question that narrows it to a kind of move and before any arrow;
\item name the habit, then its replacement: say what kind of move he played and how (its scan step, the seconds it got, what it left behind, as the digest's \texttt{your move} and \texttt{it allowed} lines say), and show how the segment's rule card replaces that habit at the board;
\item and then what: after the better move, play or say their best reply from the digest's engine line and why the move still holds, so the line stops where its gain is clear;
\item pitch it at his level: say whether this is a miss most players at his rating make or one he could have found, from how often players at his band play the best move there (the \texttt{find} the Maia step wrote on that turning point in tools/coach-review/.cache/analysis.json, or \texttt{npm run --silent coach -- engine find <fen> <uci> --rapid <his rating>}), in words and never as a percentage (under about 15\%: a pattern to learn; half or more: his to find); on a hung piece, the mirror instead, as the video grammar says;
\item praise with a board: on a move he got right, name what it saw.
\end{itemize}
\end{quote}
}
\newcommand{\AppCommit}{3a377486e8}
\newcommand{\AppCommitDays}{13}
\newcommand{\LaterCommit}{3ff95c4b72}
\newcommand{\LaterCommitDays}{six}
\newcommand{\LaterCommitList}{Oct 12, Oct 13, Oct 14, Oct 15, Oct 17 and Oct 19}
\newcommand{\PinMovePR}{136}
\newcommand{\PortPR}{161}
\newcommand{\AllDays}{23}
\newcommand{\Cases}{19}
\newcommand{\CasesCap}{19}
\newcommand{\UnusableN}{four}
\newcommand{\Unusable}{September 26, October 8, October 10 and October 11}
\newcommand{\FirstDay}{September 25}
\newcommand{\LastDay}{October 19}
\newcommand{\TrainN}{ten}
\newcommand{\HeldN}{nine}
\newcommand{\TrainCases}{Sep 25, Sep 27, Sep 29, Oct 1, Oct 4, Oct 6, Oct 12, Oct 14, Oct 17 and Oct 19}
\newcommand{\HeldCases}{Sep 28, Sep 30, Oct 2, Oct 5, Oct 7, Oct 9, Oct 13, Oct 15 and Oct 18}
\newcommand{\GamesMin}{one}
\newcommand{\GamesMax}{14}
\newcommand{\WriterRuns}{44}
\newcommand{\NullN}{six}
\newcommand{\NullDays}{Sep 25, Sep 27, Sep 28, Sep 29, Sep 30 and Oct 2}
\newcommand{\WordsPerSecond}{2.6}
\newcommand{\MinDecoyWords}{40}
\newcommand{\Spent}{205}
\newcommand{\Cap}{240}
\newcommand{\PilotCap}{60}
\newcommand{\WriterSpend}{125}
\newcommand{\GraderSpend}{79}
\newcommand{\SevenDayFirst}{39}
\newcommand{\SevenDayLast}{53}
\newcommand{\PublishedN}{19}
\newcommand{\PubOpenQuestion}{6}
\newcommand{\PubHabitNamed}{80}
\newcommand{\PubReplacement}{97}
\newcommand{\PubThenWhat}{20}
\newcommand{\PubLevel}{0}
\newcommand{\PubLevelTrue}{0}
\newcommand{\PubPraiseNamed}{69}
\newcommand{\PubOpenQuestionAll}{31}
\newcommand{\PubMoments}{3.9}
\newcommand{\PubScenes}{10.1}
\newcommand{\BaseOpenQuestion}{8}
\newcommand{\BaseHabitNamed}{87}
\newcommand{\BaseReplacement}{92}
\newcommand{\BaseThenWhat}{21}
\newcommand{\BaseLevel}{0}
\newcommand{\BaseLevelTrue}{0}
\newcommand{\BasePraiseNamed}{81}
\newcommand{\BaseOpenQuestionAll}{28}
\newcommand{\BaseMoments}{4.1}
\newcommand{\BaseScenes}{12.5}
\newcommand{\CandOpenQuestion}{74}
\newcommand{\CandHabitNamed}{69}
\newcommand{\CandReplacement}{65}
\newcommand{\CandThenWhat}{74}
\newcommand{\CandLevel}{57}
\newcommand{\CandLevelTrue}{91}
\newcommand{\CandPraiseNamed}{86}
\newcommand{\CandOpenQuestionAll}{56}
\newcommand{\CandMoments}{6.3}
\newcommand{\CandScenes}{14.7}
\newcommand{\PubSplit}{0.0}
\newcommand{\PubPositions}{3.9}
\newcommand{\PubToldHabitNamed}{80}
\newcommand{\PubToldReplacement}{97}
\newcommand{\PubToldLevel}{0}
\newcommand{\BaseSplit}{0.2}
\newcommand{\BasePositions}{4.1}
\newcommand{\BaseToldHabitNamed}{87}
\newcommand{\BaseToldReplacement}{92}
\newcommand{\BaseToldLevel}{0}
\newcommand{\CandSplit}{2.2}
\newcommand{\CandPositions}{4.3}
\newcommand{\CandToldHabitNamed}{89}
\newcommand{\CandToldReplacement}{85}
\newcommand{\CandToldLevel}{83}
\newcommand{\MovesTable}{%
\begin{tabular}{lrrr}
\toprule
In the moment scenes & Published & Today's prompt & Candidate \\
\midrule
Asks what he would play before any hint & 6\% {\footnotesize (65)} & 8\% {\footnotesize (77)} & 74\% {\footnotesize (82)} \\
Names the habit behind the mistake & 80\% {\footnotesize (74)} & 87\% {\footnotesize (77)} & 69\% {\footnotesize (120)} \\
Ends on the habit that replaces it & 97\% {\footnotesize (74)} & 92\% {\footnotesize (77)} & 65\% {\footnotesize (120)} \\
Plays out their best reply to the better move & 20\% {\footnotesize (74)} & 21\% {\footnotesize (72)} & 74\% {\footnotesize (77)} \\
Says whether the miss is normal at his level & 0\% {\footnotesize (74)} & 0\% {\footnotesize (77)} & 57\% {\footnotesize (120)} \\
\quad and that agrees with the evidence & -- & -- & 91\% {\footnotesize (68)} \\
\midrule
\multicolumn{4}{l}{\emph{Counted once per position taught}} \\
Names the habit behind the mistake & 80\% {\footnotesize (74)} & 87\% {\footnotesize (77)} & 89\% {\footnotesize (81)} \\
Ends on the habit that replaces it & 97\% {\footnotesize (74)} & 92\% {\footnotesize (77)} & 85\% {\footnotesize (81)} \\
Says whether the miss is normal at his level & 0\% {\footnotesize (74)} & 0\% {\footnotesize (77)} & 83\% {\footnotesize (81)} \\
\midrule
Praise names what the good move saw & 69\% {\footnotesize (26)} & 81\% {\footnotesize (31)} & 86\% {\footnotesize (29)} \\
Any board scene: an open question first & 31\% {\footnotesize (127)} & 28\% {\footnotesize (185)} & 56\% {\footnotesize (185)} \\
\bottomrule
\end{tabular}
}
\newcommand{\ShareBase}{43}
\newcommand{\ShareCand}{69}
\newcommand{\ShareDiff}{+0.26 {\footnotesize (+0.18 to +0.34)}}
\newcommand{\BaseHeard}{9.0}
\newcommand{\BaseWords}{1192}
\newcommand{\BaseCost}{2.61}
\newcommand{\BaseMinutes}{8}
\newcommand{\BaseRounds}{2.5}
\newcommand{\BasePassed}{19}
\newcommand{\BaseRuns}{19}
\newcommand{\BaseBoardScenes}{12.5}
\newcommand{\BaseQuestions}{9.7}
\newcommand{\CandHeard}{10.9}
\newcommand{\CandWords}{1448}
\newcommand{\CandCost}{3.17}
\newcommand{\CandMinutes}{10}
\newcommand{\CandRounds}{3.1}
\newcommand{\CandPassed}{19}
\newcommand{\CandRuns}{19}
\newcommand{\CandBoardScenes}{14.7}
\newcommand{\CandQuestions}{15.0}
\newcommand{\PubHeard}{5.9}
\newcommand{\LengthCap}{1.3}
\newcommand{\LengthRatio}{1.23}
\newcommand{\LengthRatioInterval}{1.23 {\footnotesize (1.16 to 1.29)}}
\newcommand{\TrainTruth}{+0.10 {\footnotesize ($-$0.20 to +0.40)}}
\newcommand{\TrainCoaching}{+0.80 {\footnotesize (+0.50 to +1.00)}}
\newcommand{\TrainListening}{$-$0.30 {\footnotesize ($-$0.60 to 0.00)}}
\newcommand{\TrainOverall}{+0.40 {\footnotesize (0.00 to +0.80)}}
\newcommand{\HeldTruth}{$-$0.11 {\footnotesize ($-$0.44 to +0.22)}}
\newcommand{\HeldCoaching}{+0.89 {\footnotesize (+0.67 to +1.00)}}
\newcommand{\HeldListening}{$-$0.56 {\footnotesize ($-$0.89 to $-$0.22)}}
\newcommand{\HeldOverall}{+0.78 {\footnotesize (+0.33 to +1.00)}}
\newcommand{\AllTruth}{0.00 {\footnotesize ($-$0.26 to +0.26)}}
\newcommand{\AllCoaching}{+0.84 {\footnotesize (+0.68 to +1.00)}}
\newcommand{\AllListening}{$-$0.42 {\footnotesize ($-$0.63 to $-$0.21)}}
\newcommand{\AllOverall}{+0.58 {\footnotesize (+0.26 to +0.90)}}
\newcommand{\AllNullCoaching}{+0.17 {\footnotesize ($-$0.33 to +0.67)}}
\newcommand{\TruthWins}{three}
\newcommand{\TruthLosses}{three}
\newcommand{\TruthTies}{13}
\newcommand{\CoachingWins}{16}
\newcommand{\CoachingLosses}{zero}
\newcommand{\CoachingTies}{three}
\newcommand{\ListeningWins}{zero}
\newcommand{\ListeningLosses}{eight}
\newcommand{\ListeningTies}{eleven}
\newcommand{\OverallWins}{13}
\newcommand{\OverallLosses}{two}
\newcommand{\OverallTies}{four}
\newcommand{\NullDecisive}{three}
\newcommand{\Decided}{168}
\newcommand{\SecondShare}{60}
\newcommand{\CoachingSplit}{three}
\newcommand{\SameN}{19}
\newcommand{\SameBroken}{zero}
\newcommand{\DecoyN}{19}
\newcommand{\DecoyCoaching}{zero}
\newcommand{\DecoyLostListening}{19}
\newcommand{\DecoyLostOverall}{19}
\newcommand{\DecoyAdded}{265}
\newcommand{\BaseFalse}{0.47}
\newcommand{\BaseUnsupported}{1.8}
\newcommand{\BaseMisleading}{2.7}
\newcommand{\BaseBoardFalse}{0.79}
\newcommand{\CandFalse}{0.53}
\newcommand{\CandUnsupported}{2.0}
\newcommand{\CandMisleading}{2.6}
\newcommand{\CandBoardFalse}{1.11}
\newcommand{\PubFalse}{1.0}
\newcommand{\PubBoardFalse}{1.3}
\newcommand{\FalseDiff}{+0.05 {\footnotesize ($-$0.32 to +0.47)}}
\newcommand{\BoardDiff}{+0.32 {\footnotesize ($-$0.21 to +0.90)}}
\newcommand{\RetestN}{six}
\newcommand{\RetestShare}{6}
\newcommand{\RetestMax}{20}
\newcommand{\FirstGradedDay}{September 25}
\newcommand{\FirstCoaching}{preferred the candidate}
\newcommand{\RuleTable}{%
\begin{tabular}{lrl}
\toprule
Condition & Estimate (95\% interval) & Holds \\
\midrule
Coaching on the train days above 0 & +0.80 {\footnotesize (+0.50 to +1.00)} & yes \\
Coaching on the held-out days above 0 & +0.89 {\footnotesize (+0.67 to +1.00)} & yes \\
Overall not below 0 & +0.58 {\footnotesize (+0.26 to +0.90)} & yes \\
False claims per script not up & +0.05 {\footnotesize ($-$0.32 to +0.47)} & yes \\
False board claims per script not up & +0.32 {\footnotesize ($-$0.21 to +0.90)} & yes \\
The check passes as often & 19 of 19 {\footnotesize (today's: 19)} & yes \\
Controls clean (noise floor: coaching) & +0.17 {\footnotesize ($-$0.33 to +0.67)} & yes \\
Coaching moves made more often & +0.26 {\footnotesize (+0.18 to +0.34)} & yes \\
At most 1.3 times as long & 1.23 {\footnotesize (1.16 to 1.29)} & yes \\
\bottomrule
\end{tabular}
}
\newcommand{\Kept}{yes}
\newcommand{\DayTable}{%
\begin{tabular}{llrccccrrr}
\toprule
 & & & \multicolumn{4}{c}{Judge, from the candidate's side} & Length & \multicolumn{2}{c}{Coaching share} \\
\cmidrule(lr){4-7}\cmidrule(lr){9-10}
Day & Split & Games & Truth & Coaching & Listening & Overall & ratio & Today & Candidate \\
\midrule
Sep 25 & train & 5 & 0 & 0 & 0 & 0 & 1.10 & 56 & 61 \\
Sep 27 & train & 4 & 0 & + & 0 & + & 1.13 & 40 & 71 \\
Sep 28 & held out & 3 & 0 & + & 0 & + & 1.59 & 27 & 69 \\
Sep 29 & train & 7 & 0 & + & $-$ & + & 1.27 & 45 & 54 \\
Sep 30 & held out & 6 & $-$ & + & 0 & + & 1.05 & 47 & 67 \\
Oct 1 & train & 3 & 0 & + & 0 & 0 & 1.34 & 50 & 74 \\
Oct 2 & held out & 4 & 0 & + & $-$ & + & 1.31 & 53 & 68 \\
Oct 4 & train & 4 & + & + & $-$ & + & 1.44 & 44 & 66 \\
Oct 5 & held out & 9 & 0 & + & $-$ & + & 1.14 & 23 & 37 \\
Oct 6 & train & 13 & 0 & + & 0 & + & 1.03 & 44 & 55 \\
Oct 7 & held out & 1 & 0 & + & $-$ & + & 1.19 & 60 & 56 \\
Oct 9 & held out & 4 & + & + & $-$ & + & 1.28 & 47 & 65 \\
Oct 12 & train & 4 & 0 & 0 & 0 & 0 & 1.07 & 40 & 77 \\
Oct 13 & held out & 14 & 0 & + & $-$ & + & 1.49 & 36 & 66 \\
Oct 14 & train & 6 & 0 & + & 0 & 0 & 1.18 & 33 & 67 \\
Oct 15 & held out & 8 & 0 & + & 0 & + & 1.21 & 40 & 100 \\
Oct 17 & train & 5 & + & + & 0 & + & 1.23 & 40 & 63 \\
Oct 18 & held out & 1 & $-$ & 0 & 0 & $-$ & 0.99 & 30 & 100 \\
Oct 19 & train & 2 & $-$ & + & $-$ & $-$ & 1.23 & 53 & 88 \\
\bottomrule
\end{tabular}
}


\title{Does the daily video coach better when its prompt asks for a coach's moves?\\[0.3em]
\large One bullet added to the daily routine's prompt, replayed on \Cases{} past days}
\author{Chess Coach research lab}
\date{}

\begin{document}
\maketitle

\begin{abstract}
The app's daily routine writes a narrated video of the player's games every day. A titled
coach's video review of a player near his level coaches with five moves: ask his move first, name the
habit behind the mistake and the habit that replaces it, play the better move on to the opponent's
best reply, say whether the miss is normal at his level, and praise a good move by what it saw. The
routine's videos make three of these rarely or never. We added one bullet asking for the five moves
to the routine's prompt and replayed the video step on \Cases{} past days, once under each prompt,
with the routine's own model and checker. Judging blind in both orders, a model preferred the
candidate's coaching on \CoachingWins{} days and today's prompt's on none, on the train days and the
held-out days alike, and preferred the candidate overall on \OverallWins{} days against
\OverallLosses. An audit against the evidence found the candidate asking his move first in
\CandOpenQuestion\% of its moment scenes (today's prompt: \BaseOpenQuestion\%), playing the better
move on to their reply in \CandThenWhat\% (\BaseThenWhat\%) and placing the miss at his level in
\CandToldLevel\% of the positions it teaches (today's prompt: \BaseToldLevel\%), with no more false
claims. The cost is length: its videos are \LengthRatio{} times as long heard once, and the judge
found them harder to follow at speaking pace on \ListeningLosses{} days. Every condition of the
rule fixed before the first run held, and the routine's prompt took the bullet.
\end{abstract}

\section{The question}

The app's daily routine has a model review the player's games of the day and write, besides the
written review, a narrated video script: scenes with a board from one of his games, a chart or
tiles, and beats with a spoken line, a caption, the moves played on the board and a pause where he
is meant to think. The app renders it with boards, captions and a synthetic voice.

A titled player's video review of the rapid games of an adult rated about 1100, \emph{Anyone Can Be
1200 In Chess. Watch This.} (\url{https://www.youtube.com/watch?v=PcGDehRSqTM}; the lab read a
summary of its captions, not the video itself), coaches with five moves:
\begin{enumerate}[nosep]
\item ask before telling: the student says what he would play before any hint;
\item name the habit behind the mistake, then the habit that replaces it;
\item ``and then what?'': follow the better move to the opponent's best reply;
\item pitch each miss at the player's level: a miss most players at his rating make is a pattern to
learn, one they find is his to find;
\item praise a good move by what it saw.
\end{enumerate}
The app's video grammar already asks for a pause to predict in every moment, findability or mirror
framing, and a rule card at the end of each moment, but the routine's prompt says only how a video
binds its boards and numbers. An audit of the \PublishedN{} published scripts of the days below
(Table~\ref{tab:moves}) found the second move in most moment scenes and the first, third and fourth
rarely. The candidate adds one bullet of five sub-points to step~4 of the daily prompt, right after
the bullet on how the video binds its boards:

\CandidateBullet

Nothing else in the prompt changes. The question is whether the video coaches better with it,
without becoming less true or much longer.

\section{The cases and the writer}

A case is a past day whose review landed on the app's main branch with a video script: \AllDays{}
days from \FirstDay{} to \LastDay. Each is a checkout of the app with the coach data as main held
it just before that day's review landed and the day's archived engine analysis; the app's own
\texttt{coach daily prepare} wrote the day's brief again, offline. The day's annotations, coach log,
briefing and note replies are put back as they landed, so only the video script is missing, and a
writer run writes the video for the review the day published. A case is usable only when the
published video script passes today's checker with those files; \Unusable{} do not, which leaves
\Cases{} days with \GamesMin{} to \GamesMax{} games each. Most cases run the app at
\texttt{\AppCommit} (\AppCommitDays{} days); the lab's pin moved while the cases were being built,
so \LaterCommitList{} run at \texttt{\LaterCommit}, which adds the lesson lines' naming of pins
(chess PR \PinMovePR) and changes outside the routine. The prompt and the video checker are the
same at both, and each day's two prompts ran on the same checkout. The split was fixed before any
run, interleaved by date: train is \TrainCases; held out is \HeldCases.

The writer runs as the routine does: Opus 5.5, headless, with nothing inherited from the session
that launched it, and a short preamble that names the day, says which steps ran and asks for the
video script alone, done when the routine's check passes. Each prompt wrote every day once
(\WriterRuns{} runs in all), and today's prompt wrote \NullDays{} a second time, to measure the
judge's noise floor.

\section{The graders}

A script with no model reads each run: whether the final check passes, the check rounds, files
changed besides the video script, commands it should not have run, the script's shape, its length
heard once (spoken words at \WordsPerSecond{} a second plus every pause), cost and minutes.

Sonnet 5.5 audits each script alone against the evidence it was written from: the brief without its
bulk, the clock's turning points, a line per game, the engine digest of the day, how often players
at his level play the engine's best move at each turning point, and for every board scene the
position as the archive replays it, what was played there and next, the engine's best move and
reply, how players at his band play there (Maia-2), the engine's plan lines, the threats, what each
piece hits and guards, what each move the scene plays does, and after each move the game did not
play, the engine's best reply. It lists the claims the evidence contradicts (\emph{false}), does not
carry (\emph{unsupported}) or carries but slants (\emph{misleading}), and for each board scene its
kind (a moment, praise, a judgement call, a plan, other), its false claims about the board, and
which of the five moves it makes.

Sonnet 5.5 also judges each pair blind, in both orders, on four criteria: truth, coaching, listening
(would he follow it once, at speaking pace) and overall. A side wins a criterion only when both
orders agree. Two controls check the judge on every day: today's script against itself, and against
a decoy that adds as many spoken words as the candidate added that day (at least \MinDecoyWords), in beats that
restate their scene with no square, move, piece, number or question the scene does not already
have.

The graders' prompts were settled on published scripts alone, before either prompt's runs were
graded. The pilot (the first three train days) then showed two faults, and both were fixed before
anything past the pilot was graded: the evidence did not carry what the writer's tools show (the
Maia-2 reading, the engine's reply to a move the game did not play, the plan lines), so the audit
graded claims the writer had read from them as unsupported; and one audit filed most of the lines
it checked and found true as misleading, so the claim list gained a grade for a line that holds,
which counts for nothing. Every script was graded again with the final graders. Graded with the
first evidence, the judge \FirstCoaching{} on coaching on \FirstGradedDay; the regrade's verdicts
are the ones reported.

\section{The rule}

A day scores $+1$ when the candidate wins a criterion, $-1$ when it loses and $0$ for a tie;
intervals are 95\% percentile bootstrap intervals over days. The candidate replaces the live prompt
only when every condition of Table~\ref{tab:rule} holds, fixed before its first run.

\begin{table}[h]
\centering
\small\RuleTable
\caption{The rule, fixed before the candidate's first run. Coaching and overall are the judge's
scores per day; the claim rows and the coaching moves are the audit's numbers per script, candidate
minus today's prompt; the length is the candidate's script heard once over today's on the same day.}
\label{tab:rule}
\end{table}

\section{What the bullet changed}

\begin{table}[h]
\centering
\small\MovesTable
\caption{The coaching moves the audit finds, pooled over the \Cases{} days: the share of the scenes
that owe a move and make it, with the number owing it in brackets. A moment scene owes the habit,
its replacement and the level; it owes the open question only when it asks him something, and the
reply only when it shows a better move. In the middle rows, consecutive board scenes on the same
position count once, and a move counts as made when any of them makes it.}
\label{tab:moves}
\end{table}

The published scripts and today's prompt name the habit behind a mistake and end on its replacement
in most moment scenes, but they rarely ask his move before giving a hint (\PubOpenQuestion\% and
\BaseOpenQuestion\%), rarely play the better move on to the opponent's best reply (\PubThenWhat\%
and \BaseThenWhat\%), and never say whether a miss is normal at his level. The candidate does all
three in most moment scenes (\CandOpenQuestion\%, \CandThenWhat\% and \CandLevel\%), and the level
it states agrees with the findability evidence in \CandLevelTrue\% of the scenes that state one.
Over the five moves, the share the audit finds in the moment scenes rose from \ShareBase\% to
\ShareCand\% (\ShareDiff{} a script; rule~7).

Scene by scene, the candidate names the habit and its replacement less often (\CandHabitNamed\% and
\CandReplacement\%). That is how it asks first: it often shows a position in two scenes, one that
asks and one that answers (\CandSplit{} board scenes a script repeat the position of the board scene
before them, against \BaseSplit{} under today's prompt), and the audit owes the habit, its
replacement and the level on each moment scene, so the scene that asks counts as a miss. Counted
once per position taught, the candidate names the habit as often as today's prompt
(\CandToldHabitNamed\% against \BaseToldHabitNamed\%) and ends on its replacement a little less
often (\CandToldReplacement\% against \BaseToldReplacement\%). This count was made after the result
and decides nothing.

\section{What the judge preferred}

\begin{table}[h]
\centering
\small
\begin{tabular}{lrrr}
\toprule
Criterion & Train (\TrainN{} days) & Held out (\HeldN{} days) & All \\
\midrule
Truth & \TrainTruth & \HeldTruth & \AllTruth \\
Coaching & \TrainCoaching & \HeldCoaching & \AllCoaching \\
Listening & \TrainListening & \HeldListening & \AllListening \\
Overall & \TrainOverall & \HeldOverall & \AllOverall \\
\bottomrule
\end{tabular}
\caption{The judge's mean score per day from the candidate's side ($+1$ a win, $-1$ a loss, $0$ a
tie), with its 95\% bootstrap interval over days.}
\label{tab:judge}
\end{table}

The judge preferred the candidate's coaching on \CoachingWins{} of the \Cases{} days, tied on
\CoachingTies{} and never preferred today's prompt's, and it preferred the candidate overall on
\OverallWins{} days against \OverallLosses{} (Table~\ref{tab:judge}). Truth was even: \TruthWins{}
wins, \TruthLosses{} losses and \TruthTies{} ties. Listening went the other way: the candidate never
won it and lost it on \ListeningLosses{} days. Its videos ask more and run longer, and the judge
found them harder to follow once at speaking pace.

Two runs of today's prompt on the same day, judged against each other, split on coaching on
\NullDecisive{} of \NullN{} days, each way at least once, for a mean of \AllNullCoaching. The judge
does pick a side between two scripts of one prompt, but not one way on every day it decides, as it
did for the candidate. Read in one order only, the judge favoured the side it read second in \SecondShare\% of
its \Decided{} decided verdicts, so a verdict that only one order gives is a tie.

\section{The controls}

Today's script judged against itself tied on every criterion on all \SameN{} days. The decoy,
today's script with on average \DecoyAdded{} spoken words added that only restate their scene, never
won coaching and lost listening and overall on all \DecoyN{} days. The judge does not
reward words that say nothing new, so the candidate's wins come from what its extra words say.

\section{Truth}

The audit found about as many false claims in the candidate's scripts as in today's (\CandFalse{} and
\BaseFalse{} a script; the difference is \FalseDiff) and a few more false claims about the board
(\CandBoardFalse{} against \BaseBoardFalse{} a script; \BoardDiff), within the noise; the candidate
shows more board scenes (\CandBoardScenes{} against \BaseBoardScenes{} a script). Unsupported and
misleading lines were about as common under both prompts (\CandUnsupported{} and
\BaseUnsupported{} unsupported, \CandMisleading{} and \BaseMisleading{} misleading a script).

\section{Length, cost and the check}

Every run of both prompts passed the routine's check (\CandPassed{} and \BasePassed{} of
\Cases). Heard once, the candidate's videos last \CandHeard{} minutes against \BaseHeard{}, with
\CandWords{} spoken words against \BaseWords{} and \CandQuestions{} questions against
\BaseQuestions; on the same day it is \LengthRatioInterval{} times as long (rule~8 allows
\LengthCap). The
writer took \CandRounds{} check rounds against \BaseRounds, \CandMinutes{} minutes against
\BaseMinutes{} and \$\CandCost{} against \$\BaseCost{} a day at list price.

\section{The graders measured twice}

The audit is noisy itself. The \RetestN{} pilot scripts audited twice with the same evidence (the
second time with a claim grade for a line that holds, which asks nothing new of the scene tallies)
moved by \RetestShare{} points of coaching share on average and by up to \RetestMax, so one day's
share in Table~\ref{tab:days} says little, while the pooled difference (\ShareDiff) is well beyond
that noise.

\begin{table}[h]
\centering
\scriptsize\DayTable
\caption{Each day: the judge's verdict from the candidate's side ($+$ a win, $-$ a loss, $0$ a
tie), the candidate's length over today's heard once, and the share of the coaching moves the audit
finds in the moment scenes.}
\label{tab:days}
\end{table}

\section{What the app takes}

The candidate passed every condition of the rule (Table~\ref{tab:rule}), so the daily routine's
prompt took the bullet as it was tested (chess PR \PortPR). Listening is its cost: the videos are
longer, ask more, and were harder to follow once on \ListeningLosses{} days. A next round could ask
for the same moves in fewer words, for instance one scene per moment that asks, waits and answers.

\section{Limits}

\begin{itemize}[nosep]
\item The five moves are one change: the result shows they help together, not which of them does.
\item A replay writes the video alone, after the day's review is in place; the routine writes it at
the end of a long day's work. The published videos are shorter (\PubHeard{} minutes heard once) than
either prompt's replays, so the length and its listening cost may differ live. Both prompts faced the
same setting on every day.
\item The graders are models. The judge picks a side between two runs of one prompt on some days and
favours the side it reads second; the audit's tallies move between two readings. The controls, the
noise floor and both orders bound these, and the graders are a different model from the writer.
\item The level a script states is checked against the same findability the writer reads (Maia-2 at
his band), so the audit checks that the writer read it right, not that Maia-2 knows his level.
\item The cases ran at two app commits; the prompt and the checker are the same at both, and each
day's two prompts ran on the same checkout.
\item The eval spent \$\Spent{} at list price (\$\WriterSpend{} on the writer, \$\GraderSpend{} on
the graders), on the owner's subscription, whose seven-day window it moved from \SevenDayFirst\% to
\SevenDayLast\%.
\end{itemize}

\end{document}
